%%%%%%%%%%%%%%%%%%%%%%%%%%%%%%%%%%%%%%%
% Cover Letter - MAGNIT TECH, Руководитель направления AI
%%%%%%%%%%%%%%%%%%%%%%%%%%%%%%%%%%%%%%%

\documentclass[]{cover}
\usepackage{fancyhdr}

\pagestyle{fancy}
\fancyhf{}

\rfoot{Page \thepage \hspace{0pt}}
\thispagestyle{empty}
\renewcommand{\headrulewidth}{0pt}
\begin{document}

%%%%%%%%%%%%%%%%%%%%%%%%%%%%%%%%%%%%%%
%     TITLE NAME
%%%%%%%%%%%%%%%%%%%%%%%%%%%%%%%%%%%%%%
\namesection{}{\Huge{Алексей Орлов}}{  \href{mailto:lex.orlov78@yandex.ru}{lex.orlov78@yandex.ru} | +7~(901)~797-54-88 |  \urlstyle{same}\href{https://www.linkedin.com/in/alex-orlov-v/}{LinkedIn}
}

%%%%%%%%%%%%%%%%%%%%%%%%%%%%%%%%%%%%%%
%     MAIN COVER LETTER CONTENT
%%%%%%%%%%%%%%%%%%%%%%%%%%%%%%%%%%%%%%

\currentdate{\today}
\lettercontent{Уважаемая команда Лаборатории ИИ,}

\lettercontent{Вижу, что вы строите AI-направление практически с нуля -- от команды до архитектуры. У меня похожий, но личный опыт: я самостоятельно спроектировал и развернул AI-агента-шлюза Jarvis на LangChain с собственным MCP-сервером, интегрированным с несколькими внешними системами. Хочу быть откровенным сразу: это личный/прикладной уровень, а не production enterprise-масштаб, который вы, судя по всему, ищете.}

\lettercontent{То, что я точно могу принести с первого дня -- это способность построить процесс и команду там, где их ещё нет:}

{\raggedright\fontspec[Path = OpenFonts/fonts/raleway/]{Raleway-Medium}\fontsize{11pt}{13pt}\selectfont
\begin{itemize}
    \item \textbf{Построение процессов с нуля:} спроектировал процесс управления изменениями и статусную модель задач для нескольких ИТ-направлений ПочтаТех
    \item \textbf{Построение и развитие команды:} провёл ассесмент 25+ сотрудников и разработал треки развития; запустил 2 внутренние школы, обучил 120+ человек
    \item \textbf{Технический бэкграунд:} к.т.н. по специальности «САПР», ранний опыт руководства разработкой (Oracle BI, ETL, PHP/MySQL) -- умею предметно разговаривать с инженерами
    \item \textbf{Метрики и бизнес-эффект:} настроил систему KPI и квартальное планирование с CTO, добившись сокращения Time-to-Market на 20\%
\end{itemize}\par}
\vspace{6pt}

\lettercontent{Готов честно сказать, где заканчивается мой опыт: я не разворачивал LLM/RAG-решения в production, не работал с Kubernetes и highload-инфраструктурой. Но я умею быстро учиться на стыке технологий и процесса, и мой личный проект с MCP-сервером показывает, что это не абстрактный интерес, а практика.}

\lettercontent{Если вам важнее человек, способный выстроить процесс, нанять и развить команду вокруг сильных технических специалистов, -- готов обсудить, как это может сработать.}

\lettercontent{Буду рад ответить на вопросы.}

\begin{flushright}
\closing{С уважением,}

\signature{Алексей Орлов}
\end{flushright}
\end{document}
